% ═══════════════════════════════════════════════════════════════════
%  Pages liminaires : dédicace, remerciements, résumé, abstract, glossaire
% ═══════════════════════════════════════════════════════════════════

% ── Dédicace ─────────────────────────────────────────────────────
\pageLiminaire{Dédicace}

\vspace{1.5cm}
\begin{center}
\begin{minipage}{0.82\textwidth}
\itshape
\large

Je dédie ce modeste travail :

\vspace{0.7cm}

À mes parents, pour leurs sacrifices, leur patience et la confiance qu'ils
n'ont jamais cessé de m'accorder.

\vspace{0.5cm}

À ma famille et à mes amis, pour leur soutien constant et leurs
encouragements tout au long de ces années.

\vspace{0.5cm}

À tous les professeurs de \nouvelleEtablissement, qui ont donné sans compter
de leur temps et de leur savoir pour nous former.

\vspace{0.5cm}

À Madame Fati, gérante du centre \yani, pour m'avoir ouvert les portes de son
institut, pour le temps qu'elle a consacré à répondre à mes questions et pour
la confiance qu'elle a placée dans ce projet.

\vspace{0.5cm}

À toutes celles et ceux qui, de près ou de loin, ont participé à
l'élaboration de ce travail.

\vspace{0.5cm}

Et à toute personne qui prendra ce rapport entre ses mains.

\end{minipage}
\end{center}

% ── Remerciements ────────────────────────────────────────────────
\pageLiminaire{Remerciements}

C'est avec une grande satisfaction que je tiens à remercier toutes les
personnes qui m'ont aidé, accompagné et soutenu jusqu'à l'aboutissement de ce
projet de fin d'études.

\vspace{0.3cm}

J'adresse en premier lieu mes remerciements à l'ensemble du corps professoral
de \nouvelleEtablissement, dont les enseignements constituent le socle sur
lequel ce travail a pu être bâti. Les notions de génie logiciel, de bases de
données, de sécurité et de conduite de projet mobilisées dans les pages qui
suivent y ont toutes été acquises.

\vspace{0.3cm}

Je remercie tout particulièrement \nouvelleEncadrantPedagogique, mon
encadrant pédagogique, pour son suivi régulier, ses relectures exigeantes et
les conseils qui ont orienté ce projet aux moments où il en avait le plus
besoin.

\vspace{0.3cm}

Je remercie également \nouvelleEncadrantEntreprise, ainsi que Madame Fati,
gérante du centre \yani, pour sa disponibilité et pour la qualité des échanges
que nous avons eus. C'est en observant le fonctionnement réel de son institut
--- le carnet de rendez-vous, les appels téléphoniques, la carte de fidélité
tamponnée --- que les besoins de cette application ont pu être formulés
autrement que de façon théorique. Les allers-retours réguliers avec elle ont
directement façonné les arbitrages de conception exposés dans ce rapport :
c'est à sa demande, par exemple, que la durée d'une prestation a cessé de
commander la grille des créneaux.

\vspace{0.3cm}

Mes remerciements vont enfin aux membres du jury, qui me font l'honneur
d'accepter d'évaluer ce travail, ainsi qu'à mes parents, dont le soutien
indéfectible ne s'est jamais démenti.

\vspace{0.5cm}

Merci à toutes et à tous.

% ── Résumé ───────────────────────────────────────────────────────
\pageLiminaire{Résumé}

Ce rapport de projet de fin d'études présente la conception et la réalisation
d'une solution logicielle complète pour \yani, un centre de beauté situé à
Casablanca. Le projet répond à une situation courante dans les petites
structures de services : une activité entièrement pilotée à la main --- carnet
de rendez-vous papier, prises de rendez-vous par téléphone et par messagerie
instantanée, vente de produits sans suivi de stock, fidélité gérée par une
carte tamponnée --- avec les conséquences qui en découlent : double
réservation d'un même créneau, absence d'historique exploitable, temps
consacré au téléphone au détriment des clientes présentes.

\vspace{0.3cm}

La solution développée est un mono-dépôt réunissant trois applications
adossées à une même base de données PostgreSQL. Une \textbf{API REST} bâtie
avec NestJS et Prisma porte l'ensemble des règles métier. Une
\textbf{application mobile} développée avec React Native et Expo permet aux
clientes de consulter les catalogues, de réserver une prestation, de commander
des produits et de suivre leur compte de fidélité. Un \textbf{back-office} web
en React donne à la gérante et à son personnel la maîtrise des commandes, des
rendez-vous, du catalogue, des horaires d'ouverture et du programme de
fidélité.

\vspace{0.3cm}

Le rapport détaille l'ensemble du cycle de développement : étude de
l'existant, expression des besoins, modélisation UML, choix techniques,
architecture en couches, implémentation et validation. Une méthode agile
inspirée de SCRUM, organisée en six sprints, a permis de livrer par
incréments et de confronter chaque version au fonctionnement réel de
l'institut. Une attention particulière a été portée aux propriétés que les
tests unitaires classiques ne savent pas démontrer : la sérialisation des
réservations concurrentes par verrou consultatif PostgreSQL, l'atomicité du
crédit de fidélité, l'idempotence de l'émission des bons de récompense et
l'anonymisation des comptes supprimés, exigée par la loi 09-08 sur la
protection des données personnelles. Vingt-cinq suites de tests, dont huit
écrivent dans une véritable base PostgreSQL, sont rejouées à chaque
modification par une chaîne d'intégration continue.

\vspace{0.3cm}

Le résultat est une plate-forme fonctionnelle, trilingue (français, arabe,
anglais), documentée et prête au déploiement, qui supprime le carnet papier,
rend la réservation possible à toute heure et donne à l'institut une visibilité
qu'il n'avait pas sur son activité.

\vspace{0.6cm}

\textbf{Mots clés :} application mobile, API REST, NestJS, Prisma, PostgreSQL,
React Native, Expo, React, TypeScript, réservation en ligne, programme de
fidélité, back-office, concurrence, sécurité applicative, internationalisation,
intégration continue.

% ── Abstract ─────────────────────────────────────────────────────
\pageLiminaire{Abstract}

\selectlanguage{english}

This final year project report presents the design and implementation of a
complete software solution for \textit{Yani Concept by Fati}, a beauty centre
located in Casablanca. The project addresses a situation common to small
service businesses: an activity managed entirely by hand --- a paper
appointment book, bookings taken over the phone and through instant messaging,
product sales with no stock tracking, and a loyalty programme run on a stamped
card --- with the consequences that follow: the same slot booked twice, no
usable history, and time spent on the phone at the expense of the customers
already in the salon.

\vspace{0.3cm}

The delivered solution is a monorepo bringing together three applications
backed by a single PostgreSQL database. A \textbf{REST API} built with NestJS
and Prisma carries all the business rules. A \textbf{mobile application} built
with React Native and Expo lets customers browse the catalogues, book a
treatment, order products and follow their loyalty account. A \textbf{web
back-office} written in React gives the manager and her staff control over
orders, appointments, the catalogue, opening hours and the loyalty programme.

\vspace{0.3cm}

The report covers the full development cycle: study of the existing process,
requirements elicitation, UML modelling, technical choices, layered
architecture, implementation and validation. An agile method inspired by SCRUM,
organised into six sprints, made it possible to deliver in increments and to
confront each version with the real workings of the salon. Particular care was
given to properties that ordinary unit tests cannot demonstrate: serialisation
of concurrent bookings through a PostgreSQL advisory lock, atomicity of loyalty
crediting, idempotent issuance of reward vouchers, and anonymisation of deleted
accounts as required by Moroccan law 09-08 on personal data protection.
Twenty-five test suites, eight of which write to a real PostgreSQL database,
are replayed on every change by a continuous integration pipeline.

\vspace{0.3cm}

The result is a working, trilingual (French, Arabic, English), documented
platform, ready for deployment, which removes the paper book, makes booking
possible at any hour, and gives the salon a visibility over its own activity
that it previously lacked.

\vspace{0.6cm}

\textbf{Keywords:} mobile application, REST API, NestJS, Prisma, PostgreSQL,
React Native, Expo, React, TypeScript, online booking, loyalty programme,
back-office, concurrency, application security, internationalisation,
continuous integration.

\selectlanguage{french}

% ── Glossaire ────────────────────────────────────────────────────
\pageLiminaire{Glossaire}

\begin{center}
\begin{tabularx}{\textwidth}{L{2.6cm} X}
\toprule
\entete{Sigle} & \entete{Signification} \\
\midrule
API      & \textit{Application Programming Interface} --- interface de programmation applicative \\
Argon2   & Fonction de dérivation de clé, lauréate du \textit{Password Hashing Competition} \\
CI       & \textit{Continuous Integration} --- intégration continue \\
CNDP     & Commission Nationale de contrôle de la protection des Données à caractère Personnel \\
COD      & \textit{Cash On Delivery} --- paiement à la remise ou à la livraison \\
CORS     & \textit{Cross-Origin Resource Sharing} --- partage de ressources entre origines \\
CSP      & \textit{Content Security Policy} --- politique de sécurité du contenu \\
DTO      & \textit{Data Transfer Object} --- objet de transfert de données \\
EAS      & \textit{Expo Application Services} --- chaîne de compilation native d'Expo \\
HSTS     & \textit{HTTP Strict Transport Security} \\
HTTP(S)  & \textit{HyperText Transfer Protocol (Secure)} \\
IHM      & Interface Homme--Machine \\
i18n     & \textit{Internationalization} --- internationalisation \\
IANA     & \textit{Internet Assigned Numbers Authority} --- source de la base des fuseaux horaires \\
JSON     & \textit{JavaScript Object Notation} \\
JWT      & \textit{JSON Web Token} --- jeton d'authentification signé \\
ORM      & \textit{Object-Relational Mapping} --- correspondance objet-relationnel \\
PFE      & Projet de Fin d'Études \\
REST     & \textit{REpresentational State Transfer} --- style d'architecture d'API \\
RGPD     & Règlement Général sur la Protection des Données \\
SGBD     & Système de Gestion de Bases de Données \\
SMTP     & \textit{Simple Mail Transfer Protocol} \\
SPA      & \textit{Single Page Application} --- application web mono-page \\
SQL      & \textit{Structured Query Language} \\
TDD      & \textit{Test-Driven Development} --- développement piloté par les tests \\
TTL      & \textit{Time To Live} --- durée de validité \\
UI / UX  & \textit{User Interface} / \textit{User eXperience} \\
UML      & \textit{Unified Modeling Language} --- langage de modélisation unifié \\
UTC      & \textit{Coordinated Universal Time} --- temps universel coordonné \\
UUID     & \textit{Universally Unique Identifier} \\
XSS      & \textit{Cross-Site Scripting} \\
\bottomrule
\end{tabularx}
\end{center}
